\section{Implementation}
\label{sec:impl}

extC is a compiler for a small systems language, written in 75{,}239 lines of C across 50
source files, that compiles to portable C11. The checker accounts for 32{,}064 of those lines
(43\%), the back end for 23{,}709, the parser and desugarer for 6{,}738, and the region planner
for 5{,}007; the two largest files are the code generator (16{,}017 lines) and the reclamation
runtime that it emits as text (4{,}212 lines).
No part of the runtime is a library: the generated C contains the runtime functions it needs,
so a binary links libc and nothing else, and a reader can audit one function of the produced
program without leaving the file.  The compiler is the only place where a policy decision is
made; the emitted program contains no configuration logic beyond preprocessor constants.

\subsection{What the compiler decides}
\label{sec:impl-decide}

Three analyses run before code generation, and each of them can remove work from the emitted
program rather than speed it up.

\paragraph{Placement.}
The escape analysis of \S\ref{sec:placement} assigns every allocation site a region level and
seeds the per-frame root table with the slots that can hold a reference.  The planner then
merges the candidate sites of a possibly cyclic store into the shallowest of them, which is what
makes the root table of the surviving region complete: a cycle can only be built inside one
region, so one root table sees all of its edges.

\paragraph{Unit tier.}
The tier of a translation unit is decided by a chain of three predicates, each of which can
short-circuit the next.  If no region in the unit can need reclamation the unit is compiled at
\emph{T0} and no reclamation code is emitted at all.  Otherwise, if the unit satisfies the
acyclicity conditions (i)--(iv) of \S\ref{sec:judgments}, it is compiled at \emph{T1}: counting
is on, sweeping is off, and a reference count that reaches zero frees the object immediately.
If neither predicate holds, the unit is compiled at \emph{T2}, where counting is off and the
region is reclaimed by incremental mark and sweep.  The decision is a property of the unit, not
of a call site, so a program cannot pay for a mechanism it does not use.

\paragraph{Per-function idleness.}
Independently of the unit tier, the idle proof of \S\ref{sec:tiers} is evaluated per function.
A function that cannot discard a reference without a fresh allocation in the same iteration
emits no reclamation code even when the unit keeps the runtime: it neither colours nor counts,
and its allocations are plain slot hand-outs.  The predicated is applied at the granularity of
functions because that is the granularity at which the proof is stated.

The consequence reaches the artifact.  In \tab\ref{tab:codesize} a shape compiled with
reclamation (\texttt{EXTC\_RECLAIM=1}) and the same shape compiled without it (\texttt{EXTC\_RECLAIM=0})
differ by an order of magnitude in generated code, and the difference is entirely runtime
functions plus the call sites that invoke them; the program text itself is unchanged.

\begin{table}[t]
\centering
\small
\caption{Generated C for the four benchmark shapes, with reclamation enabled and disabled.
The two configurations differ in the runtime and in the call sites that reach it; program
functions are the same in both.}
\label{tab:codesize}
\begin{tabular}{@{}lrrrrrr@{}}
\toprule
& \multicolumn{2}{c}{reclamation on} & \multicolumn{2}{c}{reclamation off} &
\multicolumn{2}{c}{ratio} \\
\cmidrule(lr){2-3} \cmidrule(lr){4-5} \cmidrule(lr){6-7}
shape & lines & bytes & lines & bytes & lines & bytes \\
\midrule
promote    & 4{,}370 & 278{,}838 & 745 & 27{,}747 & 5{,}87$\times$ & 10{,}05$\times$ \\
sizeratio  & 4{,}462 & 651{,}430 & 831 & 31{,}657 & 5{,}37$\times$ & 20{,}58$\times$ \\
chase      & 4{,}391 & 280{,}207 & 765 & 28{,}726 & 5{,}74$\times$ & 9{,}75$\times$ \\
cycleratio & 4{,}407 & 303{,}861 & 773 & 28{,}645 & 5{,}70$\times$ & 10{,}61$\times$ \\
\bottomrule
\end{tabular}
\end{table}

\subsection{Allocation}
\label{sec:impl-alloc}

A managed region owns slabs.  A slab is a 512\,KiB aligned block
(\texttt{EXTC\_RC\_SLAB\_BITS=19}); the whole block is handed out as slots, so slots of a
size class start at a fixed phase inside the block, and a slot address maps back to its slab by
a single mask.  Objects larger than the usable remainder of a block get a block of their own with
one slot; their interior pointers are resolved through a registry rather than the mask.

Slots are sized by rounding the request up to 8 bytes with a floor of 16.  The resulting stride
is not a power of two in general, so the size class is keyed by $stride/8$ and slot index is
computed by multiplication with a per-class reciprocal and a shift, never by division.  Each of
the 32 classes keeps a chain of slabs that have a free slot and a one-entry cache of the last
hit.  The invariant is one-directional: a slab with a hole is on the chain, and a slab on the
chain may be full, so a miss costs a scan of the chain rather than a lost slab.

\subsection{Slot metadata, reuse and zeroing}
\label{sec:impl-meta}

Each slot carries three bitmap bits (occupied, marked, and registered as the target of a
cross-region edge), a two-byte in-edge count, and a two-byte type identifier that the sweep uses
to find the reference fields; the payload itself is only 8-byte aligned, which keeps the layout
deterministic and the object header absent.  A freed slot is pushed onto the class chain by
writing its successor index into the per-slot type array, so the payload is never read or
written by the allocator and a stale payload cannot corrupt a chain.  The chain index is 16-bit,
which bounds the number of slots a chain can hold.

Clearing a fresh slot is done by a fixed-length switch on $n/8$ with fall-through cases down
to the 16-byte floor, and only the default arm calls \texttt{memset}.  The jump-table form
avoids the fortified \texttt{memset} that the C library substitutes for a known-length call:
that substitution cost about 13.5 instructions per allocation in the size-ratio shape before the
switch replaced it.

The allocation body is emitted \texttt{always\_inline} at every site, which folds the size class
into an immediate when the request is a constant, at a cost of about 2.7\,KB of machine code per
site; a build-time switch replaces it with a call.  Constants that the evaluation toggles are
all \texttt{ifndef}-guarded macros in the emitted runtime, so an experiment is one \texttt{-D}
flag on the C compiler and changes nothing else.

\subsection{Barriers and the grey set}
\label{sec:impl-barrier}

Two barriers maintain the tri-colour invariant of \S\ref{sec:transitions}: \texttt{shade}
colours a value that is about to be stored, and \texttt{shadeOld} colours the value being
overwritten.  The invariant they maintain is stronger than the classical one and simpler to
reason about: every pointer value held by the mutator is black, so no white object is reachable
from a mutated slot when the sweep decides.  A unit at T1 elides the colouring of local slots
entirely, because local slots are roots of the region and cannot hide a white object.

Greys are not a stack but a 2{,}048-cell open-addressed table whose cell for an address is
$((addr\gg3)+gen) \bmod 2048$ with a generation counter; a wrap of the generation clears the
table in one \texttt{memset} of 16\,KiB.  Every transition that resets mark bits advances the
generation, so a stale entry can never be mistaken for a live grey.

\subsection{Roots without a stack scan}
\label{sec:impl-roots}

A frame does not scan its stack.  The compiler emits, per frame skeleton, a static table of the
slots that can hold a reference, and the marker walks that table.  The table describes the owning
frame only; when a store places a reference into a caller frame, the store records the target
region as owed and the owner settles the debt before its next call, so no region is ever scanned
with an incomplete table.

Coroutine and task frames register their reference fields as roots when they are created, and a
worker thread claims the thread-local arena it was spawned with before it allocates.  A
cross-region edge is registered as one bit in the target slab, which is what lets a collector seed
a region from foreign references without looking at the mutator.  When the host of a store cannot
be attributed to a slab, the value is recorded in a process-level pinned table (open addressing,
load factor at most one half, Knuth backward shift on removal) and the region that contains it
drops the whole table on teardown.

\subsection{Trigger, wall and admission}
\label{sec:impl-trigger}

Reclamation is triggered by memory pressure, not by a byte counter alone: the predicate compares
the bytes handed out against a threshold, the threshold against a wall, and the bytes against the
wall directly.  The threshold is self-tuning: an empty collection multiplies it by
\texttt{EXTC\_RC\_SCALE\_MUL} (default 4), a productive collection resets it.  This is what makes
the runtime pay for the shapes that need collection and stay out of the way of the shapes that do
not: \tab\ref{tab:ablate} shows a shape that allocates but never retains improving by a factor of
two when reclamation is switched off, and a shape that retains improving by 55\% when it is
switched on.

The wall is $M_{max}$, either the value of \texttt{EXTC\_RC\_M\_MAX} or three times the live
set.  A large request reserves room before it is admitted: the slow path of the allocator records
the maximum request seen so far ($R_{max}$) and the threshold update keeps
$thr + R_{max} \le M_{max}$ whenever the wall was set explicitly (the default wall tracks the
live set and cannot carry such a promise).  The reservation costs no instruction on the hot path,
because only a new size class or a huge block enters the slow path.

\subsection{Incremental marking}
\label{sec:impl-mark}

Marking is a sequence of bounded steps driven by the allocation path; there is no collector
thread.  A step pops at most \texttt{MARK\_BATCH\_MAX=256} objects from the worklist.  Each
collection draws a quota of work proportional to the bytes allocated since the last one (one unit
per 8 bytes, capped at $2^{20}$ units), and aborts back to the mutator when the quota or a
wall-clock deadline of 800\,\us is reached, keeping the worklist and the cursor for the next step.

An object larger than the scan chunk is scanned in segments: the cursor (object and offset) lives
in the region, not on the mark stack, and \texttt{markResume} consumes one chunk per step.  A
chunk is \texttt{EXTC\_RC\_SCAN\_CHUNK=8192} words, and within a chunk the step checks the clock
every 512 words; if the elapsed time exceeds \texttt{EXTC\_RC\_SEG\_DEADLINE\_US=40} the step
stores the cursor and returns the number of words it consumed.  This is the mechanism behind the
pause theorem of \S\ref{sec:thm-bounds}: the step is bounded by the smaller of a full chunk and the
deadline plus one check interval, rather than by the size of the object.

\subsection{Incremental sweeping}
\label{sec:impl-sweep}

The sweep walks the occupancy bitmap word by word.  For each word it clears the occupancy and
registration bits, pushes every set bit onto the class chain, and advances; the work is
proportional to the capacity of the region, not to its live set, which is exactly the term that
\thm\ref{thm:amortized} charges to the allocator.  A step stops after a fixed number of words
(\texttt{EXTC\_RC\_SWEEP\_BUDGET=8}), but when it is about to leave a slab it may spend at most
\texttt{EXTC\_RC\_SWEEP\_FINISH\_MAX=64} extra words to finish that slab, so the number of
half-swept slabs stays bounded; if the credit is exhausted the slab is parked and the cursor
resumes there, which is sound because the mark bits of a parked slab are only reset when the
parked sweep resumes.

A region whose every slot is dead is freed without sweeping, and that test is $O(1)$ per
slot rather than a second pass: the sweep maintains two per-region counters, live and marked,
which are incremented as it walks.  A verification mode recomputes the counters by hand and
compares, and it is what found the one double-counting bug in this mechanism.

\subsection{Bounded teardown: the closing queue}
\label{sec:impl-close}

Destroying a region is where an unbounded pause would otherwise hide, because a region can own
thousands of slabs and each slab can hold hundreds of children.  Region teardown therefore does
$O(1)$ work on the calling thread: it aborts any collection in flight, drops the pinned entries
of the region, sets a \emph{dead} flag, appends the region to a closing queue, and returns.
The dead flag is what makes the immediate return sound: the ownership test that every store
consults rejects a slot in a dead region, so a value read out of a region after the call cannot
resurrect it, and the slab memory stays valid until the region is actually drained.

Draining happens on allocation points and on collection settle, with a budget (eight slabs or
100\,\us, whichever comes first) and a hard deadline of 1\,ms.  For each slab the drain walks the
occupancy bitmap, decrements the in-edge counts of every child, enqueues freed children on the
free queue, and, when the slab chain of the region is empty, frees the class buckets, the mark
stack, the cursor stack and the region header itself.  A queue that grows beyond 32 pending
regions is drained inline under the same budget, so the memory held by deferred teardown is
bounded even when a program destroys regions in a loop.

\subsection{Instrumentation and its two audiences}
\label{sec:impl-instr}

The artifact reports four different populations of time, and the evaluation is careful never to
add or subtract across them.
\texttt{[pause]} counts only steps taken inside an armed window, which makes it a lower bound
on the pause the program experiences.  \texttt{[step]} counts every sub-step the mutator thread
spends inside the runtime, which is the right quantity for a pause contract but includes the time
the thread was descheduled.  \texttt{[work]} counts only bounded hot-path steps, with the cap
$256\times4=1024$ enforced (\texttt{walk\_gtcap} must be zero for the run to be reported);
unbounded drains are reported separately in \texttt{[ubwork]} and \texttt{[worst]}, so a
tail that comes from a deliberately unbounded path cannot hide inside a budgeted one.

\begin{table}[t]
\centering
\small
\caption{The constants that bound a step, and the theorem each one serves.  All are
\texttt{-D}-overridable in the emitted runtime, so \tab\ref{tab:ablate} changes one mechanism at a
time without touching the compiler.}
\label{tab:budgets}
\begin{tabular}{@{}llr@{}}
\toprule
constant & role & default \\
\midrule
\texttt{SLAB\_BITS}          & slab block size (log2) & 19 \\
\texttt{SLAB\_CHUNK\_MIN}    & minimum slots per block & 512\,KiB \\
\texttt{MARK\_BATCH\_MAX}    & objects popped per mark step & 256 \\
\texttt{QUOTA\_SHIFT\_MARK}  & bytes per work unit & 3 \\
\texttt{QUOTA\_MAX\_MARK}    & work units per collection & $2^{20}$ \\
\texttt{QUOTA\_DEADLINE\_US} & wall-clock cap on one assist & 800 \\
\texttt{SCAN\_CHUNK}         & words per segmented scan step & 8192 \\
\texttt{SEG\_CHECK\_WORDS}   & words between clock checks & 512 \\
\texttt{SEG\_DEADLINE\_US}   & clock cap inside one chunk & 40 \\
\texttt{SWEEP\_BUDGET}       & occupancy words per sweep step & 8 \\
\texttt{SWEEP\_FINISH\_MAX}  & extra words to close a slab & 64 \\
\texttt{FREE\_BUDGET\_SLABS} & slabs drained per allocation point & 8 \\
\texttt{FREE\_DEADLINE\_US}  & time cap on one drain & 100 \\
\texttt{CQ\_MAX\_REGIONS}    & pending regions before inline drain & 32 \\
\bottomrule
\end{tabular}
\end{table}

\subsection{Artifact and gates}
\label{sec:impl-artifact}

The artifact is built with \texttt{make} and needs a C compiler and nothing else; the proofs
are Lean~4 modules that \texttt{proofs/verify.sh} rebuilds and audits, and the reclamation runtime is
exercised by a 216-case C driver that runs the same suite under \texttt{-fsanitize=address} and
\texttt{-fsanitize=undefined}.  Two properties are held by gates rather than by tests: the
instruction count of the allocation path has a frozen budget in a whole-program run, and the
generated C for every benchmark shape must be byte-identical to the artifact of record, so a
change to the emitted text is never silent.
